\begin{itemframe}{متغیرهای تصادفی نشانگر}
\itm
برای تحلیل بسیاری از الگوریتم‌ها، از جمله مسئله‌ی استخدام، از متغیرهای تصادفی نشانگر
\fn{Indicator random variables}
استفاده می‌کنیم که نوعی از متغییرهای تصادفی هستند. این متغییرها ابزار مناسبی برای تبدیل بین احتمال و امید ریاضی فراهم می‌کنند.
\itm
فرض کنید یک فضای نمونه $S$ و یک رویداد $A$ داریم.
آنگاه متغیر تصادفی نشانگر نسبت به رویداد $A$، به صورت زیر تعریف می‌گردد:
$$
I\{A\} =
\begin{cases}
1 & \text{اگر A رخ دهد} \\
0 & \text{در غیر این صورت}
\end{cases}
$$

\end{itemframe}


\begin{itemframe}{متغیرهای تصادفی نشانگر}
\itm
به عنوان مثالی ساده، بیایید امید ریاضی تعداد دفعاتی که در پرتاب یک سکه‌، «شیر» ظاهر می‌شود را محاسبه کنیم.
فضای نمونه برای یک بار پرتاب سکه برابر است با
$S = \{H, T\}$.
متغیر تصادفی نشانگر $X_H$ را برای رویداد $H$ (آمدن شیر) تعریف می‌کنیم:
$$
X_H = I\{H\} = \begin{cases}
1 & \text{اگر H رخ دهد} \\
0 & \text{اگر T رخ دهد}
\end{cases}
$$
\itm
در نحوهٔ نوشتار بالا
$X_H$
یک متغییر تصادفی است که برابر با
$I\{H\}$
قرار داده شده.
$I\{H\}$
به طور خاص به تعریف متغییر تصادفی نشانگر اشاره می‌کند.
\end{itemframe}


\begin{itemframe}{متغیرهای تصادفی نشانگر}
\itm
امید ریاضی تعداد دفعات آمدن شیر در یک پرتاب، برابر با امید ریاضی $X_H$ است:
$$
E[X_H] = 1 \cdot \Pr\{H\} + 0 \cdot \Pr\{T\} = \frac{1}{2}
$$
\itm
پس امید ریاضی متغییر تصادفی $X_H$ برابر است با احتمال رویداد $H$.
\itm
به طور کلی اگر $X_A$ یک متغییر تصادفی نشانگر برای رویداد $A$ باشد، آنگاه داریم:
$$
E[X_A] = \Pr\{A\}
$$
\end{itemframe}


\begin{itemframe}{متغیرهای تصادفی نشانگر}
\itm
شاید متغیرهای نشانگر برای چنین مثال ساده‌ایی بیش از حد پیچیده به نظر برسند، ولی در تحلیل موقعیت‌هایی که شامل تکرار آزمایش‌های تصادفی هستند، بسیار سودمندند.
 \itm
برای مثال، فرض کنید می‌خواهیم امید ریاضی تعداد شیر در $n$ پرتاب سکه را محاسبه کنیم. یک روش محاسبه احتمال وقوع $0$، $1$، $2$، و ... شیر است، اما راه ساده‌تری با استفاده از متغیرهای نشانگر وجود دارد.
 \itm
فرض کنید $X_i$ متغیر نشانگر برای رخ دادن رویداد «آمدن شیر در پرتاب $i$-ام» باشد. آنگاه متغیر تصادفی $X$ که تعداد کل دفعات شیر در $n$ پرتاب را نشان می‌دهد، به صورت زیر تعریف می‌شود:
$$
X = \sum_{i=1}^{n} X_i
$$
\end{itemframe}


\begin{itemframe}{متغیرهای تصادفی نشانگر}
\itm
برای محاسبه‌ی امید ریاضی $X$، از خاصیت خطی بودن امید ریاضی استفاده می‌کنیم (امید ریاضی مجموع برابر است با مجموع امید ریاضی):
$$
E[X] = E[ \sum_{i=1}^{n} X_i ] = \sum_{i=1}^{n} E[X_i]
$$
\itm
احتمال شیر آمدن در هر پرتاب برابر $\frac{1}{2}$ است پس:
$$
E[X_i] = \frac{1}{2} \quad \Rightarrow \quad \sum_{i=1}^{n} E[X_i] = \frac{n}{2}
$$
\itm
در نتیجه:
$$
E[X] = \frac{n}{2}
$$
\end{itemframe}


\begin{itemframe-s}{متغیرهای تصادفی نشانگر}{تحلیل مسئله‌ی استخدام}
\itm
حال دوباره به مسئله‌ی استخدام بازمی‌گردیم. می‌خواهیم امید ریاضی تعداد دفعاتی که شما یک دستیار اداری جدید استخدام می‌کنید را محاسبه کنیم. همان‌طور که پیش‌تر بحث شد، فرض می‌کنیم که متقاضیان به صورت تصادفی وارد می‌شوند.
\itm
فرض کنید $X$ متغیر تصادفی‌ای باشد که برابر است با تعداد دفعاتی که یک متقاضی جدید استخدام می‌شود. اگر بخواهیم امید ریاضی $X$ را مستقیماً از محاسبه کنیم، داریم:
$$
E[X] = \sum_{x=1}^{n} x \cdot \Pr\{X = x\}
$$
\end{itemframe-s}


\begin{itemframe-s}{متغیرهای تصادفی نشانگر}{تحلیل مسئله‌ی استخدام}
\itm
محاسبه مستقیم پیچیده خواهد بود. به جای آن، از متغیرهای نشانگر استفاده می‌کنیم. ابتدا $n$ متغیر نشانگر تعریف می‌کنیم که مشخص می‌کند آیا متقاضی $i$ام استخدام شده است یا نه.
$$
X_i = I\{\text{متقاضی i استخدام می‌شود}\} =
\begin{cases}
1 & \text{اگر متقاضی i استخدام شود} \\
0 & \text{در غیر این صورت}
\end{cases}
$$

\itm
در نتیجه:
$$
X = X_1 + X_2 + \dots + X_n
$$
\itm
طبق خاصیت متغییرهای نشانگر می‌دانیم:
$$
E[X_i] = \Pr\{\text{متقاضی i استخدام شود}\}
$$
\end{itemframe-s}


\begin{itemframe-s}{متغیرهای تصادفی نشانگر}{تحلیل مسئله‌ی استخدام}
\itm
متقاضی $i$-ام زمانی استخدام می‌شود که از همه‌ی متقاضیان $1$ تا $i-1$ بهتر باشد. از آنجا که فرض کرده‌ایم ترتیب ورود متقاضیان تصادفی است، $i$ متقاضی اول به صورت تصادفی ترتیب یافته‌اند، و هر یک از آن‌ها به طور مساوی می‌تواند بهترین باشد. بنابراین، احتمال اینکه متقاضی $i$ بهترین در میان این $i$ نفر باشد برابر با $\frac{1}{i}$ است.
\itm
اکنون می‌توانیم امید ریاضی $X$ را محاسبه کنیم:
$$
E[X] = E [ \sum_{i=1}^{n} X_i ]
= \sum_{i=1}^{n} E[X_i]
= \sum_{i=1}^{n} \frac{1}{i}
= \ln n + O(1)
$$
\itm
پاسخ این جمع با استفاده از ویژگی‌های سری هارمونیک داده شده است.
%todo does it need to add series summation properties in appendix
\end{itemframe-s}


\begin{itemframe-s}{متغیرهای تصادفی نشانگر}{تحلیل مسئله‌ی استخدام}
\itm
بنابراین در الگوریتم HIRE-ASSISTANT در حالت میانگین حدود
$\ln n$
نفر را استخدام می‌کنید.
\itm
پس اگر متقاضیان به صورت تصادفی وارد شوند، هزینه‌ی میانگین استخدام برابر است با:
$$
O(c_h \ln n)
$$
بنابراین، هزینه‌ی میانگین استخدام به طور قابل‌توجهی کمتر از هزینه بدترین حالت ($O(c_h n)$) است.
\end{itemframe-s}